% ─────────────────────────────────────────────────────────────────────────────
%  conclusiones.tex  —  Conclusiones generales del documento
% ─────────────────────────────────────────────────────────────────────────────

\chapter*{Conclusiones Generales}
\addcontentsline{toc}{chapter}{Conclusiones Generales}
\markboth{Conclusiones Generales}{Conclusiones Generales}

A lo largo de este documento se analizaron los 23 patrones de diseño del catálogo GoF,
cada uno aplicado a una problemática concreta dentro de la plataforma Airbnb. A
continuación se presentan las reflexiones transversales más relevantes.

\section*{Sobre los patrones creacionales}

[Síntesis de los 5 patrones creacionales: qué tienen en común, cuándo preferir uno
sobre otro dentro de una plataforma como Airbnb.]

\section*{Sobre los patrones estructurales}

[Síntesis de los 7 patrones estructurales: cómo facilitan la integración de módulos
y servicios externos en una arquitectura de microservicios o monolítica.]

\section*{Sobre los patrones de comportamiento}

[Síntesis de los 10 patrones de comportamiento: cómo mejoran la colaboración entre
objetos y la flexibilidad ante cambios de requisitos.]

\section*{Reflexión final}

[Conclusión global: importancia de conocer y aplicar patrones de diseño como
herramienta de comunicación y de arquitectura de software en proyectos reales.]
